% Source: sources/returns/S6_04_MAPS_MANIFOLDS/sections/manifolds.tex:1461–1669
\paragraph{Theorem D.2.8 (Proved: the oriented type of all simply connected Wang members).}
Every member of D.2.7 with $n=\pm1$, arbitrary $k>0$ and
arbitrary $r\in\mathbb Z$, is orientation-preservingly
diffeomorphic to
\[
 M_{\mathrm{conic}}=
 \bigl(\mathbb {CP}^3\setminus\operatorname{int}
       (S^2\times D^4)\bigr)
 \mathop{\cup}_{S^2\times S^3}(D^3\times S^3),
\]
where the removed sphere is the conic
$[s:t]\mapsto[s^2:st:t^2:0]$ with its oriented normal framing.
For the nonzero $a\in H^2(X;\mathbb F_2)$,
\[
 \langle a^3,[X]_2\rangle=1.
\]
This also supplies an independent check of the smooth
classification asserted for order three in
\cite[Theorem 7.1 and Appendix B]{manWang26}.
\begin{proof}
We first state precisely the classification being used, including
its two-primary information. Zhubr's Theorem 3 identifies the
oriented simply connected spin six-manifolds with fixed $H_2=G$
and $b_3=0$ with the corresponding spin-bordism classes over
$K(G,2)$, modulo automorphisms of $G$.
For $G=\mathbb Z/2$ his Theorem 2.10 and the calculation in
\S3.12 give
\[
 \Omega^{\mathrm{Spin}}_6(K(\mathbb Z/2,2))=\mathbb Z/2;
\]
the generator is the image of
$(\mathbb {CP}^3,\mathbb {CP}^3\hookrightarrow
K(\mathbb Z,2))$ under reduction $\mathbb Z\to\mathbb Z/2$
\cite[Theorems 3, 2.10, pp.795,803, and \S3.12, p.807]{manZhubr75}.
The homomorphism evaluating the cube of the universal
degree-two mod-two class is nonzero on this generator, and
therefore is an isomorphism to $\mathbb F_2$.
The automorphism group of $\mathbb Z/2$ is trivial. Hence there
are precisely two such oriented smooth types and their
distinguishing invariant is $\langle a^3,[X]_2\rangle$.
In particular, the vanishing of the ordinary integral $p_1$
does not need to be mistaken for a calculation of every refined
two-primary invariant.

The conic has normal bundle $\mathcal O(4)\oplus\mathcal O(2)$.
As an oriented real rank-four bundle on $S^2$ it is trivial:
its class is detected by $w_2$, which is zero. Its framing is
unique up to homotopy since $\pi_2SO(4)=0$.
Surgery preserves simple connectivity and spin, and its integral
homology sequence, with the attaching sphere representing twice
the generator, gives $H_2=H_3=\mathbb Z/2$ and all other reduced
groups zero except $H_6=\mathbb Z$.
The mod-two hyperplane class extends across the surgery trace,
because it evaluates to zero on the conic. Bordism invariance
therefore gives cube one on $M_{\mathrm{conic}}$.
It remains to compute this cube on the entire family.

Keep the common piece $X_D$ from D.2.7. The scalar involution
$\iota(z,v)=(z,-v)$ has exactly two fixed curves there: the
double elliptic curve $E$ of the primitive $I_2$ model, and the
rational curve $C$ corresponding to a vertical edge in the
rank-two triangulation. At $I_2$ the monomial coordinates
$(v,v^2/z)$ have involution $(-x,y)$, so both normal directions
at $E$ have sign $-1$. At the other rank-one point, the induced
translation of the elliptic factor is a nonzero two-torsion
point, with no fixed points. In a rank-two triangle the weights
are two odd normal weights and one even tangential weight.
There is one vertical edge orbit and its two endpoints are the
two distinct triple points, giving the smooth curve $C$.

The involution extends over the finite logarithmic quotient,
and introduces no additional fixed points there even when $k$
is even. Indeed, an equality $\iota x=g^j x$ on its central
covering torus would imply $jn/(3k)=0$ in $\mathbb R/\mathbb Z$
after applying $\psi$. Thus $3k\mid j$, since $n=\pm1$.
The equality would then assert that translation by $e_3/2$
fixes a point on a torus, which is impossible. Away from the
central torus the base rotation gives the same conclusion.
This replaces the odd-order shortcut in Wang's order-three
proof and establishes the fixed-set statement for every $k,r$.

We give the integral line-bundle calculation on the unchanged
piece $X_D$. In multiplicative coordinates let $h_1,h_2$ be
the elliptic period and the representative of $P$, and let $R$
be the multiplier of the contraction $F$. On an auxiliary line
with coordinate $\zeta$, use
\[
 \begin{aligned}
 h(z,v,\zeta)&=(h_1z,h_2v,-z^{-1}\zeta),\\
 F(z,v,\zeta)&=(h_2^{-6}z,Rv,h_2^{-3}R^3v^6\zeta).
 \end{aligned}
\]
The two lifts commute and their alternating Chern form is
$\xi=uw+6\psi\delta$. Their semicharacter modulo two is
$u+w+uw$, with zero value on the radical spanned by $e_0,e_3$;
direct substitution shows that all four monodromies preserve it.
The symmetry is checked, not inferred from vertical curvature.
Writing $B=h_2^{-3}R^3$, the lift
\[
 J(z,v,\zeta)=(z^{-1},v^{-1},-z^{-1}\zeta)
\]
satisfies $J^2=1$, $JhJ=h^{-1}$ and $JFJ=F^{-1}$;
the last identity is exactly $B^2=h_2^{-6}R^6$.
Multiplying $J$ by $-1$ normalizes it on the rigidifying section.
Thus these are symmetric rigidified Appell--Humbert data.
The displayed semicharacter is unchanged on overlap markings,
so the normalized fibrewise identifications patch.

Compare with the prequantum model on the flat real universal
cover, whose connection potential is $\pi i\xi(x,dx)$ and
whose lattice multipliers have that same semicharacter.
It is preserved by the linear monodromies. Its fibrewise
holomorphic structure is the Appell--Humbert structure for $\xi$.
Both families are symmetric and rigidified and have the same
alternating form. Their ratio is therefore a family of flat
line bundles of order dividing two. Its square is canonically
trivial on each fibre after rigidification: the trivializing
holomorphic section is unique with value one at the origin.
Those unique sections patch over the base. Hence the comparison
line is globally two-torsion on the regular family, rather than
merely having trivial fibrewise Chern class. Since $H^2(V;\mathbb Z)$
is free by the preceding integral graph calculation, its Chern
class is zero. The actual holomorphic line consequently has
exactly the prequantum Chern class on $V$, whose values on
fibre and swept cycles are respectively one and zero.
Finally $H_2(V)\to H_2(X_D)$ is onto by D.2.7 and
$H_1(X_D)=\mathbb Z$, so restriction
$H^2(X_D;\mathbb Z)\to H^2(V;\mathbb Z)$ is injective.
The extension $\mathcal L$ below therefore has values $(1,0)$
on $(s,t)$, including the mixed sweep value.

These factors extend over the three Mumford pieces. At $I_2$
the three valuations are
\[
 (\operatorname{ord}h_1,\operatorname{ord}h_2,
   \operatorname{ord}R)=(2,1,-3).
\]
Replacing $F$ by $Fh^3$ gives line multipliers, up to units,
$-x^{-2}y$ and $y^3$. The support function in the degenerating
coordinate is $-j(j-1)$; all its slopes are even. Hence the
stabilizer character of the resulting line $\mathcal L$ is
trivial along $E$. At the rank-two point the valuations are
$(1,0,1)$ and a support function is
\[
 \varphi(a,b)=-\frac{a(a-1)}2+3b^2.
\]
Its period differences are $-a$ and $6b+3$, as required by
the line multipliers $z^{-1}$ and $q^3v^6$.
Its vertical slopes are odd, giving the nontrivial character
along $C$; the adjacent horizontal slopes differ by one, so
$\deg(\mathcal L|_C)$ is odd. At the remaining rank-one
point the same factors extend over its primitive fan.
Units contribute no winding or parity, and translation to the
section vertex changes the vertical slopes by a multiple of six.

We now check the global mod-two extension without assuming that
the degree-$3k$ transfer is invertible. D.2.7 proved
$H_2(X_D)=\mathbb Z(s,t)$ and identified
$c_1(\mathcal L)$ with $(1,0)$ on that basis.
For $n=\pm1$ both first entries of
\[
 B_{k,n,r}=\begin{pmatrix}-2n&-6r\alpha\\0&1\end{pmatrix}
\]
are even. Its cokernel is generated by $s$ of order two, and
its second relation sends $t$ to an even multiple of $s$.
The nonzero global class $a$ therefore restricts exactly to
$\rho_2c_1(\mathcal L)$ on $X_D$. There is no universal
coefficient ambiguity on $X_D$ since $H_1(X_D)=\mathbb Z$.
This is the integral attachment step that makes the argument
valid for all $k$ and every $r$.

Finally $H^2(X;\mathbb Z)=0$ implies $c_1(X)=0$, hence spin.
Write $G=\langle\iota\rangle$ and
$\tau\in H^1(BG;\mathbb F_2)$ for the generator.
The class $a$ admits an equivariant lift $\widetilde a$:
the degree-two Borel obstruction with target
$H^2(BG;H^1(X;\mathbb F_2))$ vanishes, and a fixed point
splits the map to $BG$, killing the remaining differential
into $H^3(BG;\mathbb F_2)$.
The circle containing $\iota$ acts trivially on
$\pi_1(X_D)=\mathbb Z$, and a fixed point splits the extension,
so $\pi_1((X_D)_G)=\mathbb Z\times\mathbb Z/2$.
The difference between $\widetilde a$ and the equivariant
class of $\mathcal L$ vanishes in ordinary cohomology and thus
comes from this group's degree-two cohomology, by the
universal-cover low-degree sequence. Every fixed point induces
the same order-two element $(0,1)$ in this group. A single
adjustment by a multiple of $\tau^2$ consequently gives
\[
 \widetilde a|_C=a|_C+\tau^2,
 \qquad \widetilde a|_E=a|_E+\tau\eta,
 \quad\eta\in H^1(E;\mathbb F_2).
\]
The normal bundle of each fixed curve has $w_2=0$, since
$c_1(X)=0$ and $c_1(TC),c_1(TE)$ are even.
Both normal complex directions have sign $-1$, so their
equivariant Euler class is $\tau^4$.
The equivariant Thom isomorphism after inverting $\tau$ gives
\[
 \langle a^3,[X]_2\rangle
 =\int_C\frac{(a|_C+\tau^2)^3}{\tau^4}
  +\int_E\frac{(a|_E+\tau\eta)^3}{\tau^4}
 =\int_C a|_C=1.
\]
For the second equality, squares of degree-one classes on the
oriented elliptic curve vanish, and all other unwanted terms
have the wrong ordinary degree on a curve. The stated
classification now identifies $X$ with $M_{\mathrm{conic}}$.
\end{proof}
